\section{Implementation and Model Alignment}
\label{sec:implementation}
\subsection{Earlier prototype and simulation projection}
The Mooncake prototype adds a native C++ Controller, Application/Router
services and a vLLM connector. Request salts isolate physical KV across calls;
the gate seals completed buffers, resolves tier outcomes and drains pending
copies before cleanup. Recovery materializes local HBM before Store keys.
Section~\ref{sec:physical} reports its physical validation.

MoonKVSim uses one global nanosecond clock, per-worker schedulers, memory
ledgers and shared fluid-rate resources. Forward completion makes KV readable;
transfers reserve destinations and track source lifetime, service and commitment.
Failed GPU memory returns empty after repair. The frozen
\texttt{mooncake\_prompt\_v1} projection models prompt protection, request salt,
rendezvous routing, shared ordinal gating, independent tiers, recovery copying
and cleanup. It reserves separate recovery-primary memory on the bound target.

Table~\ref{tab:alignment} identifies remaining differences. The frozen workload
has exact prompt IDs and prescribed output counts, without generated KV tensors;
therefore decode protection earns no simulated recovery credit. Coarser source
allocation may overallocate a partial page but does not pad token sequences.

\begin{table*}[t]
\centering
\caption{Alignment with the Mooncake reference and remaining model boundaries.
Conformance of metadata does not calibrate device timing.}
\label{tab:alignment}
\begin{tabular}{@{}p{.21\textwidth}p{.37\textwidth}p{.37\textwidth}@{}}
\toprule
Aspect & Mooncake prototype & Evaluated projection \\
\midrule
Tier ownership & Source-owned independent HBM/DRAM & Same admission, commit and lifetime contract \\
Chunk visibility & Completed forward buffers; shared ordinal & Same atomic visibility and ordinal barrier \\
Recovery & Local copy locking; visibility lease & Explicit copy; source pinned until load completes \\
AIMD interval & 5 s in the accepted physical pilot & 1 s in the frozen simulation \\
Application signal & Causal GBDT plus calibrated residual CPM & Prescribed ideal durations; zero-Slack ablation \\
Protected state & Prompt and sealed decode KV tensors & Exact prompt token intervals; decode replay only \\
Allocation granularity & 16-token vLLM pages & 64-token source allocation units \\
\bottomrule
\end{tabular}
\end{table*}

The native HBM free path hides recovery keys and frees slots; the Controller
does not consult Worker leases when selecting pressure victims. The local
copy is locked, but eviction between activation and load may require fallback.
The simulator excludes pinned victims throughout materialization, a stronger
retention assumption. The differential checks do not validate this race.
AegisController metadata is in-memory; integration with Store HA durability
is outside the demonstrated prototype.

\subsection{Verification and revision-3 boundary}
The earlier harness compares 260 native policy cases and six ordinal-finalizer
sequences. Overlapping regression suites pass 127 and 184 tests. Execution
audits reconstruct workflow times and resource cleanup, with source/workload
hashes. These checks do not establish concurrent data-plane equivalence.

Three new Python modules provide a serial revision-3 pool owner, foreground
budget and event adapter. They execute groups, shared references, staged
commit, generation-safe free/read ACKs, TTL/deadline cancellation and resource
permits. The 159 reference tests include 81 gain combinations checked against
an independent objective and 41 memory-demand cases, plus sharing, copy failure,
reclaim retry, stale state and burst sequences. Another 33 old-contract tests
pass. These are metadata checks, not tensor movement or interference calibration.

The new reference does not replace the frozen DES or native connector. Native
integration needs allocator lending and reservation hooks, source leases, cancellation
ACKs and shared group state. Single-group concurrency and eager closure-search
cost remain unevaluated at production scale.
